\section{实验结果}
\label{sec:experiments}

\subsection{实验设置}
FruitTreeScanner系统实现为原生iOS应用程序，使用Swift、ARKit进行点云采集、Metal进行渲染、CoreML进行基于图像的果实检测。系统支持28种中国常见水果品种，具有可配置的颜色阈值。

由于本研究时无法获得真实果园数据，实验使用模拟点云数据进行。模拟生成具有可配置果实大小、密度和遮挡级别的虚拟果园场景点云。生成真实标注以验证检测和计数准确性。

性能评估指标包括单个果实的检测精度和召回率、尺寸果实的体积估算准确性以及树木和果园规模的产量估算准确性。

\subsection{点云处理结果}
点云预处理流水线在保留果实候选区域的同时，有效降低了噪声。统计离群值移除消除了对果实几何形状影响最小的孤立点。体素网格降采样在保持空间准确性的同时将点数减少约60\%。

基于颜色的过滤在有效去除绿色背景点（叶子、树枝）的同时实现了果实颜色点的高保留率。HSV颜色空间阈值相比RGB阈值对光照变化表现出更强的鲁棒性。

自适应DBSCAN聚类成功识别了合成数据中的球形果实候选。动态$\epsilon$计算适应了不同扫描距离下的密度变化。聚类验证过滤掉了细长或不规则形状的非果实物体。

\subsection{检测性能}
模拟数据的初步结果表明，在中等遮挡条件（30-50\%果实被遮挡）下，系统检测果实候选的精度超过80\%。召回率随遮挡级别变化更为显著，从中等遮挡场景的70\%到大部分果实隐藏在冠层内的重遮挡场景的50\%以下。

多模态融合方法相比单独的2D或3D检测提高了检测稳健性。融合减少了与实际果实不对应的虚假点云簇的误报，并减少了缺乏足够深度确认的2D检测的漏检。

检测果实的体积估算准确性取决于点云完整性。具有足够点覆盖的充分暴露果实显示体积估算误差低于15\%。部分遮挡的果实由于几何不完整表现出更大的估算误差。

\subsection{局限性}
当前评估完全基于模拟数据。真实果园环境存在额外挑战，包括可变光照、复杂背景植被、果实外观在品种内和品种间的变化以及扫描过程中由风引起的运动。需要通过真实果园数据进行田间验证以建立实际性能。

遮挡校正模型依赖于关于果实分布的统计假设，这些假设可能不适用于所有冠层结构。校正准确性取决于随果实类型和生长条件变化的冠层密度和结构特征。